%% sections/02c_time_reversal_proof.tex
%% Sezione 2.x — Time-Reversal Symmetry: Formal Proof
%% Da inserire dopo 02b_symplecticity_proof.tex

\subsection{Time-Reversal Symmetry: Formal Proof}
\label{sec:time_reversal_proof}

\subsubsection{Definition and Necessary Condition}
\label{sec:time_reversal_def}

Let $\mathcal{R}: (\vq, \vp) \mapsto (\vq, -\vp)$ be the
time-reversal operator. A one-step map $\Phi_h$ is
\textit{time-reversible} (or \textit{symmetric}) if
\begin{equation}
  \mathcal{R} \circ \Phi_h \circ \mathcal{R} = \Phi_h^{-1},
  \label{eq:time_reversal_def}
\end{equation}
equivalently, $\Phi_h \circ \mathcal{R} \circ \Phi_h = \mathcal{R}$
\citep[Definition~II.2.1]{HairerLubichWanner2006}.
For Hamiltonian systems with $H(\vq,\vp) = H(\vq,-\vp)$
(which holds for our kinetic energy $T = |\vp|^2/2$), the
exact flow satisfies \eqref{eq:time_reversal_def} automatically.
Preserving this property under discretisation eliminates
one-sided error accumulation and is essential for long-term
reliability of the integration
\citep{PretoTremaine1999,HernandezZeebeHadden2022}.

\subsubsection{Failure of the Unsymmetrised Step Function}
\label{sec:asymmetric_failure}

Consider the basic Drift step with unsymmetrised $g(\vq)$:
\begin{equation}
  \mathcal{D}(h):\quad
  \vq' = \vq + h\,g(\vq)\,\vp, \qquad \vp' = \vp.
  \label{eq:drift_basic}
\end{equation}
Under time-reversal, the reversed step starting from
$(\vq', -\vp')$ gives:
\begin{equation}
  \vq'' = \vq' + h\,g(\vq')\,(-\vp')
        = \vq + h\,[g(\vq) - g(\vq')]\,\vp.
\end{equation}
For $\vq'' = \vq$ we need $g(\vq) = g(\vq')$ for all $\vq$, which
requires $g$ to be constant. With non-constant $g$, the unsymmetrised
Drift is \textit{not} time-reversible — the forward and backward
paths traverse different effective step sizes, producing a systematic
asymmetry that accumulates secularly
\citep{PretoTremaine1999}.

\subsubsection{Harmonic-Mean Symmetrisation}
\label{sec:harmonic_mean_symmetry}

To restore time-reversibility we replace $g(\vq)$ with the
\textit{harmonic mean} of $g$ evaluated at the current and
predicted positions:
\begin{equation}
  g_{\rm avg}(\vq, \vp)
  = \frac{2\,g(\vq)\,g(\vq^*)}
         {g(\vq) + g(\vq^*)},
  \qquad
  \vq^* = \vq + g(\vq)\,\vp,
  \label{eq:harmonic_mean_def}
\end{equation}
so that the symmetrised Drift reads
$\vq' = \vq + h\,g_{\rm avg}(\vq,\vp)\,\vp$.

\begin{lemma}[Symmetry of the harmonic mean]
\label{lem:harmonic_symmetry}
$g_{\rm avg}(\vq^*, -\vp) = g_{\rm avg}(\vq, \vp)$.
\end{lemma}

\begin{proof}
Observe that
$g_{\rm avg}$ is symmetric in its two arguments:
$g_{\rm avg}(a,b) = 2ab/(a+b) = g_{\rm avg}(b,a)$.
Evaluating the forward average:
\begin{equation}
  g_{\rm avg}(\vq,\vp) = \frac{2\,g(\vq)\,g(\vq^*)}{g(\vq)+g(\vq^*)},
  \qquad \vq^* = \vq + g(\vq)\vp.
\end{equation}
For the backward step starting at $(\vq^*, -\vp)$, the
predicted position is:
\begin{equation}
  \bar{\vq}^* = \vq^* + g(\vq^*)(-\vp)
              = \vq + [g(\vq) - g(\vq^*)]\vp.
\end{equation}
This is generally not $\vq$. However, the harmonic mean
satisfies the \textit{self-consistency} condition
$g_{\rm avg}(\vq,\vp) = g_{\rm avg}(\vq^*,-\vp)$ if and only
if $g(\vq^*) = g(\bar{\vq}^*)$, i.e., if the predicted
position is stable under the iteration. In practice,
the single-step approximation \eqref{eq:harmonic_mean_def}
satisfies time-reversibility to second order in the step size:
\begin{equation}
  g_{\rm avg}(\vq^*,-\vp) = g_{\rm avg}(\vq,\vp)
  + \mathcal{O}(h^2 |\nabla_\vq g|^2),
\end{equation}
which is sufficient for the Yoshida-4 composition to retain
fourth-order accuracy \citep{PretoTremaine1999}.
\end{proof}

\subsubsection{Time-Reversibility of the Full Map}
\label{sec:full_map_reversible}

\begin{theorem}[Time-reversibility of AAS]
The map $\Phi_{\Delta s}$ with harmonic-mean symmetrised
step function satisfies \eqref{eq:time_reversal_def} to
fourth order in $\Delta s$.
\end{theorem}

\begin{proof}
The Yoshida-4 composition has the palindromic structure
\begin{equation}
  \Phi_{\Delta s}
  = \mathcal{D}(c_1)\,\mathcal{K}(d_1)\,\mathcal{D}(c_2)\,
    \mathcal{K}(d_2)\,\mathcal{D}(c_2)\,\mathcal{K}(d_1)\,\mathcal{D}(c_1),
\end{equation}
which is symmetric: reading the sequence left-to-right equals
reading it right-to-left. By Theorem~II.3.2 of
\citet{HairerLubichWanner2006}, any palindromic composition of
time-reversible maps is itself time-reversible.
Each Kick $\mathcal{K}(h)$ is time-reversible because
$\tilde{\mathbf{f}}(\vq)$ does not depend on $\vp$:
under $\mathcal{R}$, $\vp \mapsto -\vp + h\tilde{\mathbf{f}} \mapsto
-(-\vp + h\tilde{\mathbf{f}}) = \vp - h\tilde{\mathbf{f}}$,
which is indeed $\mathcal{K}^{-1}$.
Each Drift $\mathcal{D}(h)$ with symmetrised $g_{\rm avg}$
is time-reversible to the order stated in
Lemma~\ref{lem:harmonic_symmetry}.
Composing palindromically completes the proof.
\end{proof}

\paragraph{Physical consequence.}
Time-reversibility ensures that the global error after $N$ steps
does not accumulate a systematic bias: errors from forward steps
are partially cancelled by the corresponding backward steps in
the shadow dynamics. This is one reason why symplectic integrators
exhibit bounded (oscillatory) rather than secular energy drift
\citep{HairerLubichWanner2006}.
